\let\negmedspace\undefined
\let\negthickspace\undefined
\documentclass[journal]{IEEEtran}
\usepackage[a5paper, margin=10mm, onecolumn]{geometry}
%\usepackage{lmodern} % Ensure lmodern is loaded for pdflatex
\usepackage{tfrupee} % Include tfrupee package

\setlength{\headheight}{1cm} % Set the height of the header box
\setlength{\headsep}{0mm}     % Set the distance between the header box and the top of the text

\usepackage{gvv-book}
\usepackage{gvv}
\usepackage{cite}
\usepackage{amsmath,amssymb,amsfonts,amsthm}
\usepackage{algorithmic}
\usepackage{graphicx}
\usepackage{textcomp}
\usepackage{xcolor}
\usepackage{txfonts}
\usepackage{listings}
\usepackage{enumitem}
\usepackage{mathtools}
\usepackage{gensymb}
\usepackage{comment}
\usepackage[breaklinks=true]{hyperref}
\usepackage{tkz-euclide}
\usepackage{listings}
% \usepackage{gvv}
\def\inputGnumericTable{}
\usepackage[latin1]{inputenc}
\usepackage{color}
\usepackage{array}
\usepackage{longtable}
\usepackage{calc}
\usepackage{multirow}
\usepackage{hhline}
\usepackage{ifthen}
\usepackage{lscape}
\usepackage{circuitikz}
\tikzstyle{block} = [rectangle, draw, fill=blue!20, 
    text width=4em, text centered, rounded corners, minimum height=3em]
\tikzstyle{sum} = [draw, fill=blue!10, circle, minimum size=1cm, node distance=1.5cm]
\tikzstyle{input} = [coordinate]
\tikzstyle{output} = [coordinate]


\begin{document}

\bibliographystyle{IEEEtran}
\vspace{3cm}

\title{4.2.19}
\author{Gemini Solution}
\maketitle
{\let\newpage\relax\maketitle}

\renewcommand{\thefigure}{\theenumi}
\renewcommand{\thetable}{\theenumi}
\setlength{\intextsep}{10pt} % Space between text and floats


\numberwithin{equation}{enumi}
\numberwithin{figure}{enumi}
\renewcommand{\thetable}{\theenumi}

\textbf{Question}:\\
Find the direction and normal vectors of each of the following lines
\begin{align}
y = -x + 2
\end{align}
\solution \\

The given line equation can be rewritten in the standard form $ax + by = c$:
\begin{align}
x + y = 2
\end{align}

This equation can be expressed in vector form.\\
Let
\begin{align}
\vec{x} = \myvec{x \\ y}
\end{align}
The line equation can be written as:
\begin{align}
\vec{n^T}  \vec{x} = c
\end{align}
where $\vec{n}$ is the normal vector to the line.

By comparing $x + y = 2$ with the vector form, we can identify:
\begin{align}
\vec{n^T} = \myvec{1 & 1}
\end{align}
\begin{align}
c = 2
\end{align}

From $\vec{n^T}$, we find that the \textbf{normal vector} is:
\begin{align}
\vec{n} = \myvec{1 \\ 1}
\end{align}

The \textbf{direction vector}, $\vec{m}$, of the line is orthogonal to the normal vector $\vec{n}$. This means their dot product is zero, i.e., $\vec{n^T}\vec{m} = 0$.

For a normal vector $\vec{n} = \myvec{a \\ b}$, a direction vector can be $\vec{m} = \myvec{b \\ -a}$.
So, for our normal vector $\vec{n} = \myvec{1 \\ 1}$, the direction vector is:
\begin{align}
\vec{m} = \myvec{1 \\ -1}
\end{align}

We can verify this with the dot product:
\begin{align}
\vec{n^T}\vec{m} = \myvec{1 & 1}\myvec{1 \\ -1} = (1)(1) + (1)(-1) = 0
\end{align}\\

The normal vector of the line is $\vec{n} = \myvec{1 \\ 1}$.\\
The direction vector of the line is $\vec{m} = \myvec{1 \\ -1}$.

\begin{figure}[H]
    \centering
    \includegraphics[width=1\linewidth]{figs/Figure_1.png}
    \label{fig:placeholder}
\end{figure}

\end{document}